\subsection{All eight toric relative inputs in the wild class-two relation}

\begin{statement}{E.95}{Proved}{Evaluated cusp cochains, the complete toric wild obstruction and every fibre}
Use the actual source of E.77, the compact filling and localization
of E.92, and the following ordered lattices:
\[
 \begin{aligned}
 V&=\langle f_1,f_2,A_0,A_1,A_2,B_0,B_1,B_2,C_1,C_2\rangle,\\
 W&=\langle A_0,A_1,A_2,B_0,B_1,B_2,C_1,C_2\rangle,
 \qquad M=\Lambda^2V.
 \end{aligned}
 \tag{E95.1}
\]
Write $T_r,J_r$ for the retained actual graded actions of
$\tau,c$. Thus $T_1|_W=7I$, and $J_1|_W$ is minus the
permutation fixing $A_0,B_0$ and interchanging the three other
pairs. Let $N_2=N/\gamma_3N$ be the intrinsic class-two quotient.
All changes of its marking below are actual continuous
automorphisms; the full action and every quotient map are
transported along them.

There are explicit $c$-equivariant framed-cusp cocycles
$k_d$, representing the actual generalized-Jacobian points
$z_d$ of E.37. On the actual $G_K$, for $K=\mathbb Q(\zeta_3)$,
they are
\[
 k_d(g)=(A_g-I)b_d+r_d\kappa_3(g)+s_d\kappa_\zeta(g),
 \qquad \kappa_\zeta(g)=\frac{\chi(g)-1}{3}.
 \tag{E95.2}
\]
The three rational matrices $(b_d),(r_d),(s_d)$ are displayed
in the proof and retained entrywise in the certificate.
Their sums in (E95.2) are integral actual cocycles. The
ordinary Kummer cocycle $\kappa_3$ uses the E.49 coherent
positive roots, so its values on $\tau,\tau'$ are $1,-1$.
Put $K=(k_d(\tau))_d$. Its exact value is
\[
4K=\begin{pmatrix}
0&-2&2&0&2&-2&2&-2\\
4&0&2&8&0&-2&4&2\\
-4&-12&10&-4&6&-8&8&-14\\
2&-28&8&4&-4&-22&2&-20\\
-6&-10&22&-12&20&2&14&-8\\
-4&-15&11&-8&9&-9&5&-13\\
-6&-6&5&-6&-2&5&-10&-5\\
2&-9&4&2&-5&-4&1&4\\
-4&-11&9&-10&8&4&0&6\\
-4&-15&7&-2&-6&-12&-12&-12
\end{pmatrix}.
 \tag{E95.3}
\]
The columns have the eight-cusp order in (E95.1).

Here is the complete wild class-two image over the entire
global abelian localization image of E.92. Its local abelian
coordinates are
\[
 ([U],A)\in\mathcal P_1\times W,
 \qquad
 \mathcal P_1=V/(I-T_1)V^{J_1}\simeq C_3^3\oplus\mathbb Z_3^6.
 \tag{E95.4}
\]
For $A=\sum_d a_de_d$, the actual obstruction is
\[
 \begin{aligned}
 \omega(U,A)&=7\sum_d a_d[U-k_d(\tau),e_d]\in M,\\
 \delta(U,A)&=[\omega(U,A)]\in
 \mathcal B_2=\frac{M}{(T_2-7I)M+(J_2+I)M}.
 \end{aligned}
 \tag{E95.5}
\]
Its free coordinate is zero. Its fifteen finite coordinates
have moduli
\[
             (m_1,\ldots,m_{15})=(3^{[11]},9^{[4]}).
 \tag{E95.6}
\]
The following formulas specify every coordinate and its
inverse using the retained exact integral Smith maps.
Let $L_2$ be the E.74 left Smith matrix of
$[T_2-7I\mid J_2+I]$, and let $S$ select its fifteen
nonunit nonzero rows. Let $R_1$ be the nine representative
columns of the inverse first-degree left Smith matrix in
E.87, after its unit coordinate has been set to zero.
For $p=(p_1,\ldots,p_9)$, with $p_1,p_2,p_3\in C_3$,
write $U=R_1p$. Define
\[
 \begin{split}
 B_d(v)&=7[v,e_d],\\
 F_d&=SL_2B_dR_1,\\
 c_d&=-SL_2B_dk_d(\tau),\\
 f(p,a)&=\sum_d a_d(c_d+F_dp)\pmod{(m_i)}.
 \end{split}
 \tag{E95.7}
\]
All eight $15\times10$ matrices $[c_d\mid F_d]$ are
evaluated in \texttt{a\_toric\_wild.json}; the formula uses
only the actual matrices and the numerical columns (E95.3).
It contains no unevaluated arithmetic cochain. The exact
image condition is the fifteen congruences
\[
                        f(p,a)=0.
 \tag{E95.8}
\]
Their constant $15\times8$ matrix is
\[
 (c_d)=\begin{pmatrix}
0&0&0&0&1&2&0&0\\
0&1&2&0&2&1&0&0\\
0&2&2&0&1&1&0&0\\
0&2&2&1&1&0&0&0\\
2&0&0&0&2&1&2&2\\
1&0&2&0&0&1&1&1\\
2&0&0&0&0&0&2&2\\
1&0&0&1&0&0&1&0\\
1&0&2&2&0&2&2&0\\
0&0&0&1&2&0&2&1\\
1&0&2&1&1&2&1&1\\
8&8&2&0&0&0&0&0\\
3&0&0&0&0&0&3&3\\
3&4&7&3&5&8&3&3\\
0&7&7&2&1&7&6&6
\end{pmatrix}.
 \tag{E95.9}
\]

Every nonempty fibre over a fixed abelian class $(p,A)$
has the complete compact coordinates
\[
 \mathcal P_2\oplus\mathcal K_2
 \simeq\mathbb Z_3^{46}\oplus C_3^9\oplus C_9^5,
 \quad
 \mathcal P_2=M/(I-T_2)M^{J_2},\quad
 \mathcal K_2=\ker((I-J_2)(T_2-7I)).
 \tag{E95.10}
\]
A single saved rational matrix $Q$ gives an integral
particular relative correction precisely on (E95.8):
\[
 y_0=Q\omega(U,A),\qquad
 (I-J_2)(T_2-7I)y_0=-(I-J_2)\omega(U,A).
 \tag{E95.11}
\]
With the saved saturated basis $K_2$ and inverse $K_2^{-1}$
of E.75, the forward map takes
\[
       u=(U,z),\qquad a=(A,y_0+K_2\eta),
       \qquad [z]\in\mathcal P_2,\quad\eta\in\mathbb Z_3^{22},
 \tag{E95.12}
\]
in the actual $J$-linear class-two coordinates constructed
below. The inverse is the complete geometric gauge
normalization followed by $[z]$ and
$K_2^{-1}(y-y_0)$. All finite carries are retained by
the saved Smith maps and their inverses.

The projection of (E95.8) to $W$ is a completely evaluated
finite congruence locus: exactly $16713$ of the $9^8$
residues modulo $9W$ occur. Its full two-digit list,
the row-reduction maps, and reconstructed horizontal
solutions are retained in \texttt{a\_toric\_residues.json}.
The following table also gives every horizontal kernel:
\[
\begin{array}{c|r|r|c}
\text{first digit}&\text{number}&
 \substack{\text{second digits}\\\text{per first digit}}
 &\substack{\text{horizontal}\\\text{zero-locus type}}\\\hline
0&1&675&\mathbb Z_3^6\oplus C_3^3\\
\text{nonzero axes of }W^{J_1}\bmod3&6&486&\mathbb Z_3^6\oplus C_3^3\\
\text{other admissible digits}&54&243&\mathbb Z_3^6
\end{array}
 \tag{E95.13}
\]
An axis here means one of the three lines generated by
$A_1-A_2,B_1-B_2,C_1-C_2$. The table concerns admissible
second digits in each row; all other residues are excluded.
Consequently the complete section fibre over an admissible
fixed $A$ is, noncanonically and with explicit Smith
coordinate maps,
\[
 \mathbb Z_3^{52}\times C_3^{9+t(A)}\times C_9^5,
 \qquad t(A)=3\text{ in the first two rows of (E95.13), and }0
 \text{ in the third}.
 \tag{E95.14}
\]

For all nonunit relative values there is a short closed
form. Write $A=3a$ and put $\ell=p_6+p_8\pmod3$. Then
(E95.8) is exactly
\[
       (C+\ell D)\bar a=0\quad\text{in }\mathbb F_3^4,
 \tag{E95.15}
\]
where
\[
 C=\begin{pmatrix}
2&2&2&0&0&0&0&0\\0&0&0&0&0&0&0&0\\
0&1&1&0&2&2&0&0\\0&1&1&2&1&1&0&0
\end{pmatrix},\quad
 D=\begin{pmatrix}
1&1&1&0&0&0&0&0\\0&0&0&0&0&0&1&1\\
0&2&2&0&2&2&0&0\\0&2&2&2&1&1&0&0
\end{pmatrix}.
 \tag{E95.16}
\]
Each of the three kernels has dimension five, and their
pairwise intersections are the same three-dimensional
$W^{J_1}/3W^{J_1}$. Their union has $675$ elements.
In particular all $9W$ admits every horizontal input.
More precisely, the largest submodule of $W$ with this
last property is exactly
\[
                      9W+3W^{J_1}.
 \tag{E95.17}
\]
There are also explicit exclusions with a global abelian
preimage. For $A=A_0+B_0$, the sum of rows three and four
of $f(p,a)$ is $1\pmod3$ for every $p$. Thus every global
abelian class localizing to that relative value fails to
lift to a global intrinsic class-two section. Existence
and all fibres of its global abelian preimages are the
actual localization statement E.92. The other twenty-seven
two-distinct-cusp sums are excluded as well by the retained
exact affine systems.

\emph{Manuscript: \S51.16, folios733v--736v,
scans172--178, formulas(371)--(379), Part E(c),(d),
follow-up(4). This evaluates the wild class-two obstruction
on the entire global abelian localization image. Global
good-place relations from E.73 and higher intrinsic
obstructions remain part of E.23.}
\end{statement}

\begin{proof}
\emph{The actual finite marking and its coordinate morphism.}
Let $P=A_0$, $Q=A_2$, and use the actual framed point
$z_P$ of E.24 and E.37. Its compact Kummer class has the
representative $f_2$ on $\tau$: if the $c$-fixed
$\mathcal O_3$ Tate basis $t_E$ lifts $P$, the compatible
division points give $(g-1)t_E/3$, and
$(\tau-1)t_E/3=\zeta_3t_E$. The actual triple relation is
\[
 3z_P=(3^{-2},-3^{-1},\lambda^{-1},\lambda^{-2},
                   \zeta_3^2\lambda^{-3}),
 \quad\lambda=1-\zeta_3.
 \tag{E95.18}
\]
Divide this relation in cohomology by the literal elliptic
carry $c_T$ of E.92. The free toric part of
$\kappa(z_P)-[f_2]$ is then
$(3^{-1},3^{-1},\lambda^{-1},\lambda^{-1},\lambda^{-2})$.
Its only remaining ambiguity is the three actual global
$\mu_3$ coordinates, with exponents
$(\epsilon_1,\epsilon_2,\epsilon_3)$. The sign in the
second rational factor has the coherent odd roots all
equal to $-1$, so its pro-three Kummer cocycle is zero.

The two actual functions
\[
 f_{10}=\frac{x+\zeta_3^2y+z}{x+y},\qquad
 f_{01}=\frac{\zeta_3^2x+y+z}{x+y}
 \tag{E95.19}
\]
have divisors $A_1+B_0+C_1-3O$ and
$A_0+B_1+C_2-3O$. Their vector-frame evaluations at
$z_P$ are $1/3$ and $1/(3\lambda)$, up to the literal
sign $(-1)^{\operatorname{ord}_P f+\operatorname{ord}_O f}$.
The coherent odd-root comparison for that sign is fixed.
The toric character row is the divisor row; its two
compact columns are $-(\text{divisor row})U/3$, because
the character kills the actual map $j_*=[3I;U]$.
These are integral rows. The same possible Tate boundary
$6a$ must occur on both free generators $\tau,\tau'$.
Substitution of the two actual Kummer values therefore
gives exactly
\[
       (\epsilon_1,\epsilon_2,\epsilon_3)
                   =(1-t,t,1+t),\qquad t\in\mathbb F_3.
 \tag{E95.20}
\]
The remaining vector $(-1,1,1)$ is $j(P)$ in the three
toric torsion factors, by the actual divisor Weil pairing
of E.37. The certificate also verifies it through the
same coefficient and cohomology maps.

Choose the representative $t\in\{0,1,2\}$ of this actual
phase and set $u=1+3t$. Changing the actual Tate basis
to $u t_E$ gives the following map on $V$:
\[
 P_u=I+t j_*\pi_E,\qquad
 P_uP_v=P_{uv},\qquad
 P_u^{-1}=I-\frac{t}{1+3t}j_*\pi_E.
 \tag{E95.21}
\]
It fixes $W$, commutes with $T_1,J_1$, and fixes both
unit-character rows in (E95.19). The changed cohomology
coordinates are $P_u^{-1}$. The certificate gives the
actual integral $J_1$-fixed column $h$ satisfying
\[
 P_u^{-1}k_P^{\rm old}(\tau)-k_P(\tau)=(T_1-I)h,
 \qquad
 k_P(\tau)=(0,1,-1,1/2,-3/2,-1,-3/2,1/2,-1,-1).
 \tag{E95.22}
\]
Thus the new exponents are $(1,0,1)$.

This is also an actual full-group change of marking.
For each of the two handle generators $x_i$, send
$x_i$ to $x_i^u$ followed by the ordered product of
peripheral powers $tU_{di}$, and fix each peripheral
generator. Its Frattini map is invertible, so the
resulting continuous endomorphism of the finite-rank
free pro-three group is an automorphism $\Psi_u$.
Transport the full action and all geometric maps by
\[
 \rho_{\rm new}(g)=\Psi_u^{-1}\rho_{\rm old}(g)\Psi_u,
                \qquad F_{\rm new}=F_{\rm old}\Psi_u.
 \tag{E95.23}
\]
The inverse automorphism follows from the universal
property and the inverse Frattini map; the displayed
matrix is its abelianization. Formula (E95.23), rather
than an equality of higher actions deduced from matrices,
specifies the higher comparison. The previously retained
finite joint character changes by
$\eta_{\rm new}=\eta_{\rm old}+t\pmod3$, as direct
multiplication of its E.89 character row by $P_u$ shows.
The old phase and old character are retained through
this map, without a numerical assignment to either.

\emph{All eight actual framed points.}
Use coherent roots with
$\kappa_\lambda=(\kappa_3+\kappa_\zeta)/2$, from the
literal identity $\lambda^2=-3\zeta_3$. E.92 gives the
full cochain identity
\[
 3u_{f_2}(g)=c_{T,3}\kappa_3(g)+c_{T,\zeta}\kappa_\zeta(g)
                              +(A_g-I)h_*.
 \tag{E95.24}
\]
It holds on both $\tau,\tau'$ and at $c$; both sides
are cocycles on the actual free split source, so it
holds everywhere. The exact columns $h_*,c_{T,3},c_{T,\zeta}$
are retained in the certificate. Adding the phase-zero
toric tuple gives the full $k_P$, not only its value on
one Frobenius.

The actual CM automorphism
$\sigma:[x:y:z]\mapsto[x:y:\zeta_3z]$ sends $P$ to
$Q$ and fixes $O$. If $S$ is its literal cusp permutation
and $R=\left(\begin{smallmatrix}0&-1\\1&-1\end{smallmatrix}\right)$,
its map on $V$ is
\[
 \Sigma=\begin{pmatrix}R&0\\(UR-SU)/3&S\end{pmatrix}.
 \tag{E95.25}
\]
This follows from the torus permutation, the compact CM
map, and equivariance of the specified $j$; these three
maps uniquely determine the lower block. The matrix is
integral, has cube $I$, and commutes with the actual
$G_K$ action. The vector-frame factors at $Q$ and $O$
are both $\zeta_3^{-1}$ and are $1$ elsewhere. Hence
\[
       \sigma_*z_P=z_Q+\Delta,
 \quad \Delta_Q=1,\quad\Delta_d=\zeta_3\ (d\ne O,Q),
       \qquad k_Q=\Sigma k_P-\kappa(\Delta).
 \tag{E95.26}
\]

For $a+b=d$ in $E[3]$, divide the chord or tangent
through $a,b$ by the line through $d,-d,O$; when
$d=O$, use the tangent at $O$ as denominator. The
resulting function $f$ has divisor $a+b-d-O$.
Multiplication by this actual function in the framed
line bundles proves
\[
       z_a+z_b-z_d=(\operatorname{lead}_r(f)/
                         \operatorname{lead}_O(f))_{r\ne O}.
 \tag{E95.27}
\]
All leading coefficients are computed in the literal
parameters $x/z,y/z,z/x$ at $A_i,B_i,C_i$. For a line
$L=ax+by+cz$, the first nonzero entries of the following
lists give its order and leading coefficient:
\[
\begin{array}{c|ccc}
A_i&(0,c-b\zeta_3^i)&(1,a)&(3,-b\zeta_3^i/3)\\
B_i&(0,c-a\zeta_3^i)&(1,b)&(3,-a\zeta_3^i/3)\\
C_i&(0,a-b\zeta_3^i)&(1,c)&(3,-b\zeta_3^i/3).
\end{array}
 \tag{E95.28}
\]
They follow directly by substituting the cubic equation
in the indicated local parameter. Every value is a
literal $\pm\zeta_3^e\lambda^n$. The operations
$P+P,Q+Q,P+Q,P+B_1,P+B_2,B_0+B_2$ generate all eight
points from $P,Q$. The certificate verifies all sixty-four
addition identities. In each comparison the same integral
boundary column works on $\tau$ and $\tau'$. It also
verifies every triple relation, proper Kummer projection,
and conjugation identity. Averaging the cocycle at $d$
with the $c$-transform of that at $cd$ preserves its actual
class and gives the chosen $c$-equivariant representative.

For completeness the entire matrices in (E95.2) are
as follows, with the same row and column orders as $K$:
\[
24(b_d)=\begin{pmatrix}
8&4&0&16&-4&0&4&8\\0&4&-4&0&-4&4&-4&4\\
0&-2&2&0&2&-2&2&-2\\4&0&2&8&0&-2&4&2\\
4&2&0&8&-2&0&2&4\\0&-3&3&0&3&-3&3&-3\\
2&0&1&4&0&-1&2&1\\6&3&0&12&-3&0&3&6\\
2&-3&1&4&0&-4&2&-2\\6&3&3&12&0&0&6&6
\end{pmatrix},
 \tag{E95.29}
\]
\[
6(r_d)=\begin{pmatrix}
0&0&0&0&0&0&0&0\\0&0&0&0&0&0&0&0\\
-4&-1&-1&-2&-2&-2&-3&-3\\-1&-4&-1&-2&-2&-2&-3&-3\\
-1&-1&-4&-2&-2&-2&-3&-3\\-2&-2&-2&-4&-1&-1&-3&-3\\
-2&-2&-2&-1&-4&-1&-3&-3\\-2&-2&-2&-1&-1&-4&-3&-3\\
-3&-3&-3&-3&-3&-3&-6&-3\\-3&-3&-3&-3&-3&-3&-3&-6
\end{pmatrix},
 \tag{E95.30}
\]
\[
6(s_d)=\begin{pmatrix}
0&0&0&0&0&0&0&0\\0&0&0&0&0&0&0&0\\
0&-8&8&0&5&-5&8&-8\\5&-18&7&10&-3&-16&5&-11\\
-5&-7&18&-10&16&3&11&-5\\0&-10&10&0&7&-7&7&-7\\
-2&-3&5&-1&0&4&-5&-1\\2&-5&3&1&-4&0&1&5\\
1&-5&10&-1&8&5&6&9\\-1&-10&5&1&-5&-8&-9&-6
\end{pmatrix}.
 \tag{E95.31}
\]
These rational summands have the exact symmetries
$b_{cd}=J_1b_d$, $r_{cd}=-J_1r_d$, $s_{cd}=J_1s_d$.
Their actual cocycle sum is integral, as established
by the point and frame construction. Substitution of
$\kappa_3(\tau)=1$, $\kappa_\zeta(\tau)=2$ proves
(E95.3). In particular (E95.2) replaces the framed
division input at every actual arithmetic element
on this toric cusp subspace.

\emph{Literal geometric paths and class-two coordinates.}
Start with actual $c$-equivariant peripheral paths.
For a rational cusp, coprime complement conjugacy in
its pro-three path torsor gives a fixed path; for a
conjugate pair, choose one path and its conjugate.
The framed Abel--Jacobi comparison of E.37, including
its coherent odd-root sign transport, identifies the
abelian path class with $\kappa(z_d)$. If its initial
representative is $k_d^{\rm old}$, the unique boundary
column carrying it to (E95.2) is obtained from
\[
 b_d^{\rm path}=(T_1-I)^{-1}
                 (k_d^{\rm old}(\tau)-k_d(\tau)).
 \tag{E95.32}
\]
It is integral because the two actual cocycles have
the same class. Uniqueness uses $H^0(P,V)=0$, or
equivalently the invertibility of $T_1-I$ over
$\mathbb Q_3$. The corresponding equation on $\tau'$
and the $c$-equivariance are retained. Choose group
lifts of these columns equivariantly; fixed columns
lift to fixed elements by successive coprime lifting
through the lower central quotients. Conjugating the
peripheral generators by these actual lifts performs
the required change. Its Frattini map is the identity,
so it extends to an automorphism of the full free group.
All higher actions and maps are again transported as
in (E95.23).

At class two, use the integral BCH lattice, valid
because $2$ is a unit in $\mathbb Z_3$. Choose a
linear lift of $V$, and average it with its $J$-transform:
\[
 l(v)=\tfrac12\bigl(l_0(v)+J l_0(J_1v)\bigr),
 \qquad
 (v,m)(v',m')=(v+v',m+m'+\tfrac12[v,v']).
 \tag{E95.33}
\]
The lift remains the specified peripheral lift on $W$.
It gives an actual Lie and group coordinate isomorphism,
with $J(v,m)=(J_1v,J_2m)$. It can also be lifted to a
full free-group automorphism by choosing lifts of its
ten generator images. The peripheral action in these
coordinates is the literal formula
\[
 T(e_d,0)=(7e_d,-7[k_d(\tau),e_d]).
 \tag{E95.34}
\]
It follows by expanding the actual path conjugation
of the positive peripheral generator. Linearity of
the Lie action gives its value on every $A\in W$.
The action on the other handle lifts is retained as
its actual transported action; it is not replaced by
the graded matrix when a gauge operation uses it.

\emph{The obstruction and its integral inverse.}
Normalize a local section to have value $1$ at $c$.
Write its two independent values as $u=(U,z)$ and
$a=(A,y)$. The actual source relation of E.77 says
that $uT(a)u^{-1}a^{-7}$ is fixed by $J$.
For $A\in W$, its abelian coordinate is zero and
its central coordinate, by (E95.33)--(E95.34), is
\[
                 (T_2-7I)y+\omega(U,A).
 \tag{E95.35}
\]
The powers contribute no additional BCH term because
$T_1A=7A$. Applying $I-J_2$ gives the claimed exact
relation. The factor-one-half map (E75.7) sends its
cohomology class to $[\omega]$ in $\mathcal B_2$.
Every bracket in $\omega$ contains a peripheral
vector, so the free compact elliptic coordinate
vanishes.

For $v\in V^{J_1}$ the literal identity
\[
 7[(I-T_1)v,A]=-(T_2-7I)[v,A]
 \tag{E95.36}
\]
proves descent through the complete first-degree
gauge. A change $k_d\mapsto k_d+(T_1-I)v_d$
also changes its contribution by the displayed
image of $T_2-7I$. Thus both the path morphisms
and the finite carries in $\mathcal P_1$ are
accounted for. The certificate checks them as
literal coefficient identities.

Let $H^-$ be the last block of the inverse splitting
matrix of $J_2$ in E.75, and write its exact Smith
identity as
\[
 L^- H^-(I-J_2)(T_2-7I)R^-=D^-.
 \tag{E95.37}
\]
There are twenty-three nonzero diagonal entries,
with valuations $0^{[8]},1^{[11]},2^{[4]}$.
Taking the inverse of these entries and setting
the remaining coordinates to zero gives the actual
rational matrix
\[
 Q=-R^-_{*,1:23}(D^-_{1:23,1:23})^{-1}
                   L^-_{1:23,*}H^-(I-J_2).
 \tag{E95.38}
\]
Its product with (E95.37) proves (E95.11) on the
forty-four-dimensional cusp-containing summand.
The exact diagonal divisibilities say that $Q\omega$
is integral precisely when (E95.8) holds. This proves
both directions of the image condition and supplies
an actual section whenever it holds. The saturated
kernel and its inverse give all relative solutions.

To give the complete geometric quotient, first use
the unique fixed abelian conjugator to carry $U$
to its Smith representative $R_1p$. Lift that
conjugator to the actual $J$-fixed subgroup of
$N_2$ and apply the literal operations
$u\mapsto w uT(w)^{-1}$, $a\mapsto w a w^{-1}$.
These operations use the full transported action
on $w$ and retain its central terms. The remaining
fixed conjugators are central and change only $z$
by $(I-T_2)M^{J_2}$. Their quotient has Smith
valuations $0^{[7]},1^{[9]},2^{[5]}$, with twenty-four
free coordinates. The actual inverse Smith matrices
and the saturated relative kernel inverse prove
(E95.10)--(E95.12). No stabilizer is divided out
after this step: its first nonzero graded coordinate
would be a fixed vector for $T_r$, whose eigenvalues
are the positive geometric weights of E.37, and none
is $1$. This also proves uniqueness of the aligning
conjugators.

\emph{The complete finite image and its inverse algorithm.}
Equations (E95.7) are affine linear in $p$ and linear
in $a$. Their coefficients are reduced row by row
with the actual moduli (E95.6). The three finite
$p$ coordinates can be changed by $3$ without
altering any residue. Thus the system factors through
$A\bmod9W$, with the other six horizontal residues
taken modulo $9$. For a fixed $A$, exact Smith
elimination on
\[
 [\,\sum a_dSL_2B_d\mid\operatorname{diag}(m_i)\,]
 \tag{E95.39}
\]
decides the affine consistency, gives a particular
$U$, and gives its complete integral kernel modulo
the fixed gauge. The certificate retains both
inverse matrices for forty-eight actual systems,
including every single cusp and every two-distinct-cusp
sum. Formula (E95.39) is the same finite algorithm
for arbitrary input coefficients.

The exhaustive residue calculation has a shorter
two-stage implementation whose intermediate maps
are retained. For each $a_0\in\mathbb F_3^8$, solve
the fifteen equations modulo $3$ as
$p_0=x_0+N x$. There are exactly sixty-one consistent
$a_0$. Write $a=a_0+3a_1$, use the integer lifts
of these matrices, and let $F_0,c_0$ be the four
modulo-nine rows at $a_0$. Set
\[
 t_0=(F_0x_0+c_0)/3,\qquad Q_0=F_0N/3,
 \qquad Z\text{ a basis of the left kernel of }F_0\bmod3.
 \tag{E95.40}
\]
The divisions are exact integer divisions. If
$F_1,c_1$ are the same rows evaluated on $a_1$,
the remaining condition is precisely
\[
 Z(Q_0+F_1N)x=-Z(t_0+F_1x_0+c_1)\pmod3.
 \tag{E95.41}
\]
A solution gives $p_0=x_0+Nx$, and the literal
equation
$F_0p_1=-(t_0+Q_0x+F_1p_0+c_1)\pmod3$
reconstructs $p=p_0+3p_1$. Thus the finite computation
has a proved inverse at every step. Gaussian
elimination over $\mathbb F_3$ on all $6561$ first
digits, and on all $6561$ second digits for each
consistent first digit, gives (E95.13). The saved
lists are complete, and reconstructed witnesses
are checked in the original fifteen congruences.

The kernel on the three finite horizontal coordinates
depends only on $a_0$, because their modulo-nine
coefficients are divisible by $3$. Its dimension
is exactly $t(A)$ in (E95.14). A closed subgroup
of $C_3^3\oplus\mathbb Z_3^6$ defined by finite
congruences has free rank six; its torsion is
exactly that computed kernel. Its Smith maps
give continuous forward and inverse coordinates.
Combining them with the linear particular correction
(E95.38), and then the central fibre maps, proves
the full fibre (E95.14), including its topology.

For $a_0=0$, divide the last four rows by $3$.
Their complete bilinear tensor factors as the
single scalar $p_6+p_8$, giving (E95.15)--(E95.16).
The three matrices have rank three and any two
stacked have rank five. Their common kernel is
the span of the three cusp differences fixed by
$J_1$. This proves the count $3\cdot3^5-2\cdot3^3=675$
and every nonunit assertion. Finally intersect
the kernels of all constant and all horizontal
coefficient rows, with their actual moduli. Its
exact Hermite basis is
\[
\begin{pmatrix}
9&0&0&0&0&0&0&0\\0&9&6&0&0&0&0&0\\
0&0&3&0&0&0&0&0\\0&0&0&9&0&0&0&0\\
0&0&0&0&9&6&0&0\\0&0&0&0&0&3&0&0\\
0&0&0&0&0&0&9&6\\0&0&0&0&0&0&0&3
\end{pmatrix}.
 \tag{E95.42}
\]
It is exactly the Hermite basis of $9W+3W^{J_1}$.
Both inclusions and the inverse lattice coordinates
are verified by integer multiplication. For the
explicit exclusion $A_0+B_0$, adding rows three
and four of the original system gives the constant
$1\pmod3$, with every horizontal coefficient zero.
E.92 supplies actual global abelian preimages;
functoriality of localization transfers the failure
of the local intrinsic class-two lift to each of them.
This concludes the exact wild calculation while
retaining the other global relations of E.73.
\end{proof}
